\documentclass[11pt]{article}
\usepackage[margin=0.9in]{geometry}
\usepackage{fontspec}
\usepackage{amsmath,amssymb,amsthm}
\usepackage{xurl}
\begin{document}\sloppy
\begin{center}
{\Large GROK H0012: first-step cube 1 of \texttt{a40ff54e580f5d0a} at $L=9$ closes UNSAT through all 210 subcubes, and 151 of cube 0's 210 subcubes are certified}\\[4pt]
2026-10-06 00:12 America/Los\_Angeles \quad branch \texttt{main}, tip \texttt{1d665862} \quad automation 41774ea2\\[2pt]
\textbf{P versus NP: OPEN.}
\end{center}

\section*{Peer delta since H2354 (\texttt{1d665862})}
Checked 00:13\,PT. Nothing landed on \texttt{main} after \texttt{1d665862} (06:57 UTC, 23:57\,PT). Nothing under \texttt{research/main/} or \texttt{research/agents/grok/} is new. Root \texttt{main.tex} is not touched by this note.

\section*{1. Literature}
None needed. The next claim is about our own finite search, and the only question is whether the split pieces cover the parent cube. That is answered by reading our own scripts, below.

\section*{2. Claims recorded}
\textbf{Claim CUBE1 (COMPUTATIONAL, solver-trusted).} First-step cube 1, the step $(w,s_a,s_b,g)=(0,0,1,4)$, of the tight $L=9$, $r=2$, $n=6$ search for \texttt{a40ff54e580f5d0a} is UNSAT. Phase B of \texttt{split\_queue.sh} wrote \texttt{CUBE-DONE 1 00:09:32}. Its 210 subcube lines name 210 distinct indices and 210 distinct step-2 tuples, all UNSAT under CaDiCaL, with no SAT and no TIMEOUT line for cube 1. Solve times: median 11.6\,s, slowest 305.6\,s, sum 5499.6\,s, 135 over 1\,s. These subcubes carry no DRAT certificate.

\textbf{Claim DRAT151 (COMPUTATIONAL).} 151 of cube 0's 210 subcubes carry drat-trim-verified refutations. The 8 new ones are 88, 164, 130, 132, 27, 208, 26 and 124. That leaves 59, and subcube 190 was being checked at 00:12\,PT.

\textbf{Claim HANDOFF-74 (COMPUTATIONAL, about our own pipeline).} Cube 74, the orphan of the stopped runner, logged \texttt{TIMEOUT\_or\_exit=124} by 00:04\,PT, the log's last write. So the guard's catch-up set is now 72 and 74, plus 69 if phase B missed it. The guard (\texttt{phaseB\_catchup.sh}, pid 1459005) is alive and waiting for \texttt{split\_queue.sh} (pid 1403476) to exit.

\section*{3. Smallest honest test}
\begin{itemize}
\item \textbf{Index control} (inline Python, 00:13\,PT). Phase B splits cube $k$ by indexing \texttt{regsat\_fix.cubes(6)}. The runner that timed the cube out indexes \texttt{regsat\_cube.cubes(6)}. If the two lists differed, ``cube 1 is closed'' would refer to the wrong cube. Both lists have 150 entries and are equal, entry by entry. Entry 1 is $(0,0,1,4)$, which matches the tuple printed on every phase-B cube-1 line. Indices $k<75$ are exactly the $w=0$ half that W-SWAP keeps.
\item \textbf{Shared step-2 map} (same run). For every $j$ in $0,\dots,209$, subcube $j$ of cube 1 fixes the same step-2 tuple as subcube $j$ of cube 0. So cube 1 was split by the same 210-way map whose cover argument CUBE0 already relies on, and nothing new needs checking.
\item \textbf{Recount} (\texttt{/workspace/scratch\_h2333\_1005/experiments/recount.py}, rerun 00:12\,PT). The 75 trivial-UP lines plus 6 more give 81. The 129-entry queue is disjoint from them, and together they partition $\{0,\dots,209\}$. The H2301 batch has written 70 lines in queue order, all \texttt{s VERIFIED}, with 70 distinct CNF hashes. That gives 151 certified. Solve times run 7.8 to 64.6\,s, check times 7.9 to 127.3\,s, and check/solve 0.90 to 2.58.
\item \textbf{Coverage control} (runner log, 00:13\,PT). Cubes 0 to 74 each have exactly one result line: 42 UNSAT, 33 TIMEOUT, 0 SAT. Cubes 0 and 1 are now closed by splitting, so 44 of the 75 $w=0$ cubes are closed and 31 remain open. Phase B has started cube 2 and has 82 of its subcubes UNSAT, none SAT.
\end{itemize}

\section*{4. What this changes}
One fewer open first-step cube, from 32 to 31. The bounds stay $9\le L(\texttt{a40ff54e580f5d0a},2)\le 11$ (COMPUTATIONAL). The lower bound would rise to 10 only if all 31 open cubes also close. The largest confirmed price stays one.

\section*{Effect on the standing candidate}
No claim is falsified and no charged resource improves. The time-state-representation surface stays CANDIDATE. The SAT-to-charged-model compiler is OPEN, and so is the product-span obstruction. These results cover one fixed function and say nothing about lower bounds for any family.

\section*{Next smallest step}
Let phase B work through the remaining timed-out cubes, and let the DRAT batch finish cube 0's last 59 subcubes. Once \texttt{split\_queue.sh} exits, check that the guard's \texttt{CATCHUP} lines name 72, 74 and, only if phase B skipped it, 69.

\bigskip
\noindent\textbf{Status sentence.} P versus NP remains OPEN: the SAT-to-charged-model compiler and the product-span obstruction are open load-bearing arrows. This pass records that first-step cube 1 of \texttt{a40ff54e580f5d0a} at tight $L=9$ is UNSAT through all 210 of its subcubes (COMPUTATIONAL, solver-trusted, with the cube index checked against the runner's list), that 151 of cube 0's 210 subcubes now carry drat-trim-verified refutations (COMPUTATIONAL), and that cube 74 timed out and falls to the catch-up guard. Among the 75 kept first-step cubes, 44 are closed and 31 are open, with 0 SAT, and $9\le L(f,2)\le 11$ is unchanged. No claim is falsified, the largest confirmed price stays one, no charged resource improves, and this is not a proof of anything about P versus NP.
\end{document}
